\clearpage
\section*{September 9, 2026: Checking five paired 99-unit proposals}
\textit{Agent-drafted evidence review.}
\textit{An extension below records the subsequent completion of the ten-pair check.}

Five selected post-policy parents consisting of two 99-unit filings have substantial evidence of coordinated development. The cases are 230/234 Clarkson Avenue, 937/951 East 108 Street, 1029/1035 Atlantic Avenue, 940/942 Woodycrest Avenue, and 23/29 Mount Hope Place. This is a purposive case review, not an estimate of the error rate among all parent links.

The East 108 Street, Woodycrest, and Mount Hope filings explicitly identify their companion job and state that the buildings share one zoning lot, owner, and applicant. Outside reporting describes the Clarkson and Atlantic Avenue pairs as two 99-unit buildings. A September 8 financing report also identifies the East 108 Street buildings as one 198-unit development. Woodycrest appears under a combined address in the designer's project index. These sources support retaining the original links; they do not establish final separate policy assessment.

Mount Hope requires an additional distinction. The original two filings each report 99 units, but September 1 reporting describes 99 residences in one building and 246 units across the site. All 21 records for the two job roots in the frozen local DOB history report 99 proposed units. The broader site figure could reflect additional development or a reporting problem and is not resolved here. It does not justify imputing another filing or replacing the July snapshot. Confidence in the original connection is higher than confidence that 198 exhausts the current site proposal.

No production membership, unit count, or treatment classification changed. The initial recommendation was to review the remaining exact-198 parents, followed by larger repeated-99 groups and ordinary links in both periods. Nearby unlinked filings should also be checked: validating accepted pairs cannot reveal every missed companion.

\paragraph{Extension: all ten pairs checked.} The remaining five cases were subsequently inspected. Reporting supports the Atlantic/Havens, Queens, and Grand Concourse/Sheridan pairs. Shepherd/Highland has strong administrative support, but the cited article confirms only the Shepherd building. Third Avenue has a consequential measurement difference: the Housing Database-priority panel records 99+99, while the saved initial DOB records report 98+97. The pair has related acquisition and filing evidence, but a later separate sale also qualifies continuing common ownership. These findings do not establish an erroneous link; they require resolving the intended unit measure before fitting the exact-198 moment. Detailed sources are in \path{2026-09-09-ten-pair-followup.md} in the same audit report directory. The current estimation sample contains 25 historical and 51 recent multi-constituent parents; at that stage, only the ten exact-198 cases had received this external case check. The next entry extends the review to all 76 estimation parents.

\textit{Reproduction note.} Review of code revision \texttt{8458668}, the July 10 DOB capture with filings through July 8, and linked parent/constituent panels. Case identifiers, filing dates, source URLs, access limitations, and judgments are recorded in \path{tasks/audits/audit_symmetric_parent_cohorts/report/2026-09-09-five-99-pairs.md}. External sources were consulted September 9; later reports provide retrospective corroboration and were not used as pre-filing characteristics. This entry is compiled with \texttt{make} in \path{logbook/}.
